\section{Related Work}
\label{sec:related_work}

The intersection of large language models (LLMs) and text clustering has catalyzed a shift from purely statistical grouping to semantically driven, human-aligned categorization. 
This section contextualizes our system within recent methodological precedents spanning conversational preference elicitation, LLM-guided clustering mechanics, and instruction-tuned text embeddings.

\subsection{Conversational Preference Elicitation}
Traditional preference learning often relies on passive feedback or static pairwise interactions. 
Recent studies demonstrate that multi-turn natural language dialogues significantly increase informative value, reduce user friction, and expose nuanced criteria that passive settings miss \cite{eliciting_human_preferences}. 
To optimize these dialogues, existing frameworks frequently depend on complex architectural additions or expensive post-training phases. 
For instance, TO-GATE relies on trajectory optimization to balance clarification resolvers with task-aligned responses \cite{to_gate}, while other approaches leverage diffusion-inspired funneling strategies to fine-tune LLMs for incremental profile generation \cite{asking_clarifying_questions}.

In contrast to these computationally heavy fine-tuning and optimization methods, our work introduces a streamlined, production-ready prompting paradigm. 
We exploit the strong instruction-following capabilities of Qwen3.6-35B-A3B \cite{qwen3.6} by engineering a robust system prompt that transforms the model into an adaptive interview agent. 
By feeding initial cluster insights directly into this system prompt, the agent is dynamically empowered to guide the conversation and extract actionable user preferences without requiring any underlying model fine-tuning. 
This demonstrates that sophisticated, context-aware preference elicitation can be achieved purely through strategic prompt engineering and contextual data conditioning.

\subsection{LLM-Assisted Clustering Algorithms}
Rather than treating clustering as a separate downstream task, modern frameworks inject LLM intelligence directly into the clustering loop. 
Early methods demonstrated that foundational models could act as zero- or few-shot categorizers by leveraging in-context learning to assign text fragments to candidate clusters based solely on illustrative prompt examples \cite{few_shot_clustering}. Building upon this, ClusterLLM utilizes hard triplet questions (e.g., assessing relative document similarities) and pairwise granularity queries to guide hierarchical text partitioning \cite{clusterllm}. However, traditional pairwise constraints scale poorly on extensive datasets. To alleviate this, optimized algorithms generate multi-element constraint sets (Must-Link and Cannot-Link blocks) via LLMs, achieving massive reductions in total API queries \cite{optimized_algorithms}.

Other paradigms replace numeric geometric operations entirely. 
\textit{Dial-In LLM} incorporates an LLM judge to label cluster coherence and iteratively isolate stable partitions while merging redundant configurations based on text intent labels \cite{dial_in_llm}. 
Furthermore, approaches like \textit{k-LLMmeans} replace traditional numeric vector centroids with concise, LLM-generated textual summaries, enriching semantic interpretability and coherence \cite{summaries_as_centroids}. 
In extreme cases, frameworks such as LLM-MemCluster eschew vector spaces altogether, manipulating a dynamic memory set of cluster labels with linear runtime complexity \cite{llm_memcluster}. 

\subsection{Instruction-Guided Embeddings}
While memory-only and summary-centroid systems offer distinct advantages, vector representations remain crucial for scalable text matching. 
Crucially, a fixed embedding space cannot capture subjective, task-specific groupings. 
Recent breakthroughs demonstrate that decoder-only LLMs can be transformed into generalist embedding models by replacing causal masking with bidirectional attention and integrating pooled latent layers \cite{nv_embed}. 

Models such as \textit{e5-mistral-7b-instruct} prove that embedding vectors can adapt dynamically to explicit user instructions passed directly into the encoder prompt \cite{improving_text_embeddings}. In this project, we utilize \textbf{Qwen3-Embedding (8B)} \cite{qwen3_embedding}, an advanced instruction-tuned encoder that leverages this exact paradigm. By encoding the rich task instructions derived from our Qwen3.6-35B-A3B interviewing agent, the Qwen3-Embedding model projects document vectors into a space that reflects the user's subjective criteria, modifying semantic distances dynamically without requiring downstream fine-tuning.

\subsection{Methodological Gap: Subjective Alignment}
A critical limitation within current literature is the pervasive assumption that an objective, underlying ``ground-truth'' assignment exists for every text corpus, typically validated using external metrics like Normalized Mutual Information (NMI) or Rand Index (ARI) \cite{clusterllm, few_shot_clustering, llm_memcluster, optimized_algorithms}.

Our work departs from this paradigm by addressing real-world deployments where text clustering is highly subjective and rigid mathematical benchmarks often misalign with a user's actual downstream goal. 
Because no single, universally correct cluster assignment exists for such tasks, there is a distinct scarcity of research surrounding purely user-aligned evaluation. 
Consequently, this study shifts the evaluative focus away from geometric benchmark matching. 
Instead, we formalize internal conversational metrics—such as turn efficiency, conversation length, and LLM-as-judge multi-dimensional faithfulness—to assess how effectively a system adapts its representations to a user's idiosyncratic intent.